\documentclass[12pt,a4paper]{article}
\usepackage[utf8]{inputenc}
\usepackage[portuguese]{babel}
\usepackage{amsmath}
\usepackage{amsfonts}
\usepackage{amssymb}
\usepackage[margin=2.5cm]{geometry}
\usepackage{titlesec}
\usepackage{graphicx}
\usepackage{booktabs} 
\usepackage{tabularx} \usepackage{array}
% --- Configuração para remover numeração automática das seções ---
\setcounter{secnumdepth}{0}

% --- Definição do cabeçalho ---
\pagestyle{myheadings}
\markright{}
\pretolerance=10000

\begin{document}

% --- Cabeçalho da Instituição ---
\begin{center}
    \large
    UNIVERSIDADE FEDERAL DO RIO DE JANEIRO \\
    ESCOLA POLITÉCNICA \\
    DEPARTAMENTO DE ENGENHARIA DE COMPUTAÇÃO E INFORMAÇÃO
\end{center}

\vspace{2cm}

% --- Título do Documento ---
\begin{center}
    \Large\bfseries
    PROPOSTA DE PROJETO DE GRADUAÇÃO
\end{center}


% --- Informações do Aluno e Orientador ---
\begin{center}
\noindent \textbf{Aluno:} Antonny Victor da Silva \\
\textbf{E-mail:} antonny.victor.silva@poli.ufrj.br\\
\textbf{Orientador:} Natasha Costa \\
\textbf{Curso:} Engenharia de Computação Informação.
\end{center}



\section*{\bfseries 1. TÍTULO}
Avaliação de Redes Neurais para Tomada de Decisão para Task Offloading em Ambientes Simulados de Computação Edge e Cloud

\section*{\bfseries 2. ÊNFASE}
Computação


\section*{\bfseries 3. TEMA}
O tema deste trabalho versa sobre a tomada de decisões para transferência de
tarefas computacionais (task offloading) utilizando redes neurais sem peso em infraestruturas heterogêneas de computação em borda
(Edge) e em nuvem (Cloud).


\section*{\bfseries 4. DELIMITAÇÃO}
Este trabalho delimita-se ao desenvolvimento e à comparação experimental de estratégias de task
offloading implementadas em duas camadas complementares:
1) Camada C#/.NET 10 (Simulador Analítico): desenvolvida em .NET 10 responsável pela
geração de um dataset sintético balanceado e estratificado com 15.000 amostras, pelo cálculo
analítico do tempo de resposta (Edge/Cloud) que gera os rótulos de treinamento.\\

2) Camada Python (EdgeSimPy): desenvolvida no framework de simulação por eventos discretos EdgeSimPy 1.1.0, atuando como um ambiente de validação temporal e de rede realista. Essa
camada simula de forma independente as conexões de rede de transporte, filas de servidores, e
provisionamento. O modelo de tarefas nesta camada é estendido para implementar ciclos de vida
temporais complexo.\\

O estudo limita-se ao agendamento concorrente e à avaliação estatística das decisões das políticas geradas pela camada C# quando submetidas à simulação realista de eventos discretos da camada
Python, sem abordar ainda mecanismos de mobilidade física de usuários ou migração de serviços ativos.

O projeto adota uma abordagem de extensões escalonadas: 1) Edge vs Cloud; 2) carga de rede dinâmica; 3) múltiplos servidores Edge; 4) mobilidade; 5) energia, não implementando todas as extensões simultaneamente.

\section*{\bfseries 5. JUSTIFICATIVA}
O advento de sistemas cibernéticos físicos, Internet das Coisas (IoT), veículos conectados e
algoritmos de visão computacional em tempo real impõe limites severos às arquiteturas de nuvem
centralizada. Fatores como saturação do link de rede de longa distância, flutuações severas na largura
de banda e altos tempos de propagação inviabilizam a transferência de dados pesados para servidores
em nuvem distantes. A computação em borda (Edge Computing) surge como resposta natural,
descentralizando servidores físicos e aproximando o processamento dos sensores geradores.No
entanto, os servidores de borda possuem capacidade restrita de hardware e memória. Tomar a decisão
correta sobre descarregar (offload) uma tarefa ou executá-la na borda local exige a avaliação
constante de variáveis complexas e conflitantes. Heurísticas baseadas em regras determinísticas ou
limiares simples de hardware são ineficazes em ambientes reais de flutuação dinâmica. Redes neurais
representam uma alternativa viável por sua habilidade em generalizar padrões multidimensionais nãolineares sob condições severas.Contudo, a avaliação de estratégias baseadas em inteligência
computacional frequentemente peca ao utilizar ambientes simplificados de simulação analítica. A
integração de um simulador de eventos discretos calibrado como o EdgeSimPy atua como um
laboratório dinâmico fundamental. Ele possibilita auditar o impacto das decisões de offloading neural
diretamente sobre métricas físicas simuladas, promovendo uma ponte empírica rigorosa e
reproduzível para o desenvolvimento de sistemas operacionais de borda mais eficientes.


\section*{\bfseries 6. OBJETIVO}
O objetivo central do trabalho é avaliar o desempenho sistêmico de decisões de task
offloading geradas por redes neurais artificiais (MLP e WiSARD) em um ambiente dinâmico e
restrito de simulação física de computação em borda e nuvem, priorizando métricas sistêmicas
(latência média, P95/P99, taxa de violação de deadline, throughput, utilização de recursos) sobre
métricas puras de classificação.

O projeto adota uma separação explícita de três camadas: (A) geração analítica de rótulos no
simulador C#, (B) treinamento dos modelos neurais, e (C) avaliação independente via simulação
EdgeSimPy, reconhecendo a limitação de circularidade quando A e C compartilham a mesma fórmula.

Para alcançar este objetivo geral, delimitam-se as seguintes etapas de validação:1. Treinar estaticamente os modelos neurais WiSARD
e MLP em C#, avaliando métricas clássicas de classificação (acurácia, precisão, recall e F1-score)
contra uma base rotulada sintética.2. Desenvolver um modelo temporal e hierárquico de execução de
Tasks dentro do ecossistema EdgeSimPy (Fase 5 de simulação), implementando controle estrito de
memória volátil temporária e filas de agendamento por prioridade (FIFO) por servidor.3. Integrar o
agendador de tarefas temporais ao ciclo de clock dinâmico principal do simulador (Fase 6),
permitindo que a passagem de tempo cronológico do EdgeSimPy avance de forma síncrona com os
eventos de execução de tarefas na borda.4. Acoplar a inteligência de decisão neural do C# ao motor
em Python via contratos de dados estruturados (CSV/JSON), submetendo as redes a testes de estresse
computacional e instabilidades de rede (NetworkFlows).5. Comparar sistematicamente as redes
neurais (WiSARD e MLP) contra políticas clássicas de mercado (Random, Fixed Rule baseada em
CPU e Simple Heuristic), estabelecendo um ranking empírico das métricas globais do ecossistema.

\section*{\bfseries 7. METODOLOGIA}

A metodologia científica do projeto adota uma progressão estruturada e isolada de engenharia
e simulação, dividida em fases consecutivas e documentadas:\\

\textbf{Fases Concluídas (Desenvolvimento do Motor Temporal):}
\begin{itemize}
    \item Fase 1 e 2: Extração e mapeamento dos relacionamentos físicos do dataset oficial e inicialização limpa do loop dinâmico do simulador EdgeSimPy sem gargalos.
    \item Fase 3 e 4: Desenvolvimento e comparação determinística das políticas de placement primárias (FirstFit, LatencyAware e ResourceAware), seguida por uma auditoria fina do ciclo de
provisionamento de recursos de computação na borda.
    \item  Fase 5: Criação do domínio lógico de tarefas temporárias (Task e TaskStatus) independente de serviços estáticos. Desenvolvimento do componente TaskScheduler em Python, com suporte a múltiplas tarefas, filas FIFO em cada servidor de borda, e consumo de memória temporária que é completamente liberada após a simulação temporal, validado sob rigorosos testes de asserção. Fases Atuais e Futuras (Integração e Testes Neurais):
\end{itemize}

\textbf{Fases Atuais e Futuras (Integração e Testes Neurais):}
\begin{itemize}
    \item  Fase 6 (Fase Corrente): Acoplamento dinâmico do TaskScheduler ao ciclo síncrono temporal do EdgeSimPy (Simulator.step()). Isso garantirá que o relógio físico simulado avance o processamento das tarefas no motor de eventos.
    \item Fase 7 (Transmissão de Dados): Integração da transmissão de payloads pela rede física através do agendamento dinâmico de NetworkFlows nas linhas de comunicação entre o usuário e o
servidor escolhido, permitindo modelar a latência estocástica de envio de dados.
    \item Fase 8 (Acoplamento da Inteligência): Construção da ponte de decisões em Python que invoca
os modelos WiSARD e MLP para decidir o offloading com base no estado instantâneo do
simulador. Hipótese: políticas neurais superam heurísticas simples em cenários de carga dinâmica.
    \item Fase 9 (Análise de Desempenho): Execução de simulações em processos paralelos limpos para
evitar compartilhamento indesejado de estados dinâmicos. Coleta e consolidação de métricas
globais: taxas de violação de deadline por aplicação, tempo médio de decisão (overhead), vazão
sistêmica, latência extrema (P95 e P99) e eficiência energética sob cenários estressantes de
sobrecarga computacional.
\end{itemize}

\section*{\bfseries 8. MATERIAIS}
Para a viabilização e execução do projeto de graduação, serão utilizados os seguintes recursos
computacionais:\\
\textbf{Ambiente de Inteligência e Decisão (Lado C#):}
\begin{itemize}
    \item SDK do .NET 10 para desenvolvimento em linguagem C# de alto desempenho.
    \item IDE Microsoft Visual Studio Code ou similar com extensões C# oficiais.
    \item Algoritmos de redes neurais WiSARD e MLP codificados em código puro, sem pacotes de
frameworks de aprendizado profundo (como ML.NET ou PyTorch), viabilizando portabilidade e
execução determinística na banca examinadora. 
\end{itemize}
\textbf{Ambiente de Simulação Temporal (Lado Python):}
\begin{itemize}
    \item Interpretador Python 3.12 estruturado com ambiente isolado (.venv).
    \item Simulador EdgeSimPy (versão 1.1.0, commit 76eb5ead74596bb4240759fa4336f1d6f190c70a).
    \item Bibliotecas científicas auxiliares para manipulação, processamento e modelagem física: Pandas 2.x, Numpy 1.2x, NetworkX (para caminhos mínimos e roteamento de rede via algoritmo Dijkstra), Matplotlib e Seaborn para síntese visual e plotagem dos gráficos de performance
comparativos.
\end{itemize}
\textbf{Hardware para Execução dos Experimentos:}
\begin{itemize}
    \item Computador pessoal de mesa equipado com processador multinúcleo x86-64, 16 GB de
memória RAM e sistema operacional Linux Ubuntu / Windows Subsystem for Linux (WSL) para
garantia de ambiente de compilação limpo e estável.
\end{itemize}

\section*{\bfseries 9. CRONOGRAMA}
O planejamento de execução do projeto de graduação está dividido em fases sequenciais de
engenharia, integração e redação ao longo do período letivo de 6 meses, conforme detalhado na tabela cronológica e mapeado visualmente na estrutura metodológica do projeto:\\
\begin{table}[htbp]
    \centering
    \caption{Cronograma de execução do Projeto de Graduação}
    \label{tab:cronograma}
    
    \renewcommand{\arraystretch}{1.4}
    
    \begin{tabularx}{\textwidth}{
        >{\raggedright\arraybackslash}X
        *{6}{>{\centering\arraybackslash}p{0.8cm}}
    }
        \toprule
        \textbf{Atividade / Fase de Trabalho}
        & \textbf{Mês 1}
        & \textbf{Mês 2}
        & \textbf{Mês 3}
        & \textbf{Mês 4}
        & \textbf{Mês 5}
        & \textbf{Mês 6} \\
        \midrule
        
        Fase 1: Configuração do ambiente e treinamento estático (C\#)
        & X & X & & & & \\[0.1cm]
        
        Fase 2: Desenvolvimento de Tasks e Filas no EdgeSimPy
        & & X & X & & & \\[0.1cm]
        
        Fase 3: Acoplamento temporal dinâmico (Fase 6 / NetworkFlow)
        & & & X & X & & \\[0.1cm]
        
        Fase 4: Integração de inteligência neural (C\# $\rightarrow$ Python)
        & & & & X & X & \\[0.1cm]
        
        Fase 5: Análise estatística e redação da monografia final
        & & & & & X & X \\
        
        \bottomrule
    \end{tabularx}
    
    \vspace{0.15cm}
    
    \raggedright
    \footnotesize
\end{table}

\begin{thebibliography}{99}

\bibitem{souza2023edgesimpy}
SOUZA, Paulo S.; FERRETO, Tiago; CALHEIROS, Rodrigo N.
\textit{EdgeSimPy: Python-Based Modeling and Simulation of Edge Computing Resource Management Policies}.
Future Generation Computer Systems, v. 148, p. 446--459, 2023.
DOI: \url{https://doi.org/10.1016/j.future.2023.06.013}.

\bibitem{souza2023edgesimpygithub}
SOUZA, Paulo S. et al.
\textit{edge\_sim\_py v1.1.0}.
GitHub, 2023.
Disponível em: \url{https://github.com/antonnyvictor18/EdgeCloudOffloadingTcc}.
Acesso em: 30 ago. 2026.

\bibitem{smith2005virtual}
SMITH, Jim; NAIR, Ravi.
\textit{Virtual Machines: Versatile Platforms for Systems and Processes}.
San Francisco: Morgan Kaufmann Publishers, 2005.
ISBN 1558609105.

\bibitem{souza2024alocacao}
SOUZA, Paulo S. et al.
\textit{Análise de Políticas de Alocação em Computação em Borda Usando Simulação de Eventos Discretos}.
In: SIMPÓSIO BRASILEIRO DE REDES DE COMPUTADORES E SISTEMAS DISTRIBUÍDOS (SBRC), 2024.
Anais [...]. 2024.

\bibitem{haykin2001redes}
HAYKIN, Simon.
\textit{Redes Neurais: Princípios e Prática}.
2. ed. Porto Alegre: Bookman, 2001.

\bibitem{carneiro2020wisard}
CARNEIRO, H. C. et al.
\textit{Multiclass classification with WiSARD weightless neural network}.
In: INTERNATIONAL CONFERENCE ON HYBRID INTELLIGENT SYSTEMS, 2020.
Proceedings [...]. 2020.

\end{thebibliography}


\end{document}
